%%%%%%%%%%%%%%%%%%%%%%%%%%%%%%%%%%%%%%%%%%%%%%%%%%%%%%%%%%%%
%% Lecture 15
%%%%%%%%%%%%%%%%%%%%%%%%%%%%%%%%%%%%%%%%%%%%%%%%%%%%%%%%%%%%

\lecture
{15}
{Monday, 28 September 2026}
{Introduction to Probability}
{Dr. Nishant Chandgotia}


%============================================================
% Tails and the Behaviour of S_n
%============================================================

\section{Tails and the Behaviour of Sums}

\subsection{How Does \(S_n/n\) Behave for Large \(n\)?}

Let
\(
X_1,X_2,\ldots
\)
be a sequence of random variables and define
\[
S_n=X_1+\cdots+X_n.
\]

A fundamental question is:

\begin{question}{Behaviour of the Empirical Average}
{behaviour-empirical-average}

How does
\(
\frac{S_n}{n}
\)
behave as $n\to\infty$?

\end{question}

The answer is strongly influenced by the tail behaviour of the
random variables. In particular, the quantity
\(
\mathbb{P}(|X|>t),
\qquad t\to\infty,
\)
plays an important role.

\begin{observation}

The behaviour of the tails
\[
\mathbb{P}(|X|>t)
\]
for large $t$ determines how frequently very large values of $X$
occur.

Random variables whose tails decay slowly are often referred to as
having \emph{heavy tails}.

\end{observation}


%============================================================
% A General Normalization
%============================================================

\subsection{Normalization of Partial Sums}

For suitable sequences of random variables, one may seek a sequence
of positive numbers $b_n$ such that
\(
b_n\longrightarrow\infty
\)
and
\[
\frac{S_n}{b_n}\xrightarrow{P}0.
\]

The appropriate normalization depends on the tail behaviour of the
random variables.

This motivates the study of weak laws for triangular arrays and
appropriate choices of the normalization sequence $b_n$.


%============================================================
% Fair Games
%============================================================

\section{Fair Games and Pricing}

\subsection{Pricing an Independent Game}

Suppose a casino game has winnings represented by a random variable
$X$.

If a sequence of people independently plays the same game, the
winnings give rise to an i.i.d. sequence
\[
X_1,X_2,\ldots,
\qquad
X_i\overset{d}{=}X.
\]

Let
\[
S_n=X_1+\cdots+X_n.
\]

We now ask:

\begin{question}{How Should the Casino Price the Game?}
{pricing-casino-game}

Suppose
\(
\mathbb{E}(|X|)<\infty.
\)

Then the expected winning per game is
\(
\mu=\mathbb{E}(X).
\)

If the casino charges a price $\mu'$ per game, how should $\mu'$
be chosen?

\end{question}


\subsection{Three Possible Pricing Scenarios}

Suppose
\[
\mu' > \mu.
\]

Then, on average, the casino makes money.

However, the price is higher than the expected winning and may therefore
be viewed as non-competitive.

\medskip

If
\[
\mu'<\mu,
\]
then the casino loses money on average.

\medskip

The natural fair price is therefore
\[
\boxed{
\mu'=\mu=\mathbb{E}(X).
}
\]

When
\[
\mathbb{E}(|X|)<\infty,
\]
the Weak Law of Large Numbers gives
\[
\frac{S_n}{n}
\xrightarrow{P}
\mu.
\]

Thus the total net loss of the casino after $n$ games, when the price
per game is $\mu$, is
\[
n\mu-S_n.
\]

\begin{observation}

The Weak Law tells us that
\[
\frac{n\mu-S_n}{n}
=
\mu-\frac{S_n}{n}
\xrightarrow{P}0.
\]

Thus
\[
n\mu-S_n=o_P(n).
\]

The question is whether the casino's net loss
\(
n\mu-S_n
\)
can nevertheless become large and positive.

\end{observation}


%============================================================
% Finite Variance Case
%============================================================

\subsection{What Happens When the Variance Is Finite?}

Suppose additionally that
\[
\operatorname{Var}(X)=\sigma^2<\infty.
\]

Then the Central Limit Theorem gives
\[
\frac{S_n-n\mu}{\sqrt{n}}
\xrightarrow{d}
N(0,\sigma^2).
\]

Equivalently,
\[
\frac{n\mu-S_n}{\sqrt{n}}
\xrightarrow{d}
N(0,\sigma^2).
\]

Thus, in the finite-variance case, the fluctuations of the total
winnings around their mean are typically of order
\[
\sqrt{n}.
\]

The limiting normal distribution is symmetric about zero, so the
positive and negative fluctuations are balanced asymptotically in the
sense described by the Central Limit Theorem.

\begin{question}{The Finite-Expectation Case}
{finite-expectation-case}

What happens if we only assume
\(
X\in L^1
\)
and do not assume
\(
X\in L^2?
\)

\end{question}


%============================================================
% St. Petersburg Game
%============================================================

\section{The St.\ Petersburg Game}

\subsection{Description of the Game}

Consider the following game.

A fair coin is tossed repeatedly until the first head appears.

Let
\(
N
\)
denote the number of tosses required to obtain the first head.

The winnings are defined by
\[
X=2^N.
\]

Since the first head occurs on the $k$-th toss exactly when the first
$k-1$ tosses are tails and the $k$-th toss is a head,
\[
\mathbb{P}(N=k)
=
\left(\frac12\right)^k.
\]

Consequently,
\[
\boxed{
\mathbb{P}(X=2^k)=2^{-k},
\qquad k\geq1.
}
\]

\subsection{The Expected Winning}

We have
\[
\begin{aligned}
\mathbb{E}(X)
&=
\sum_{k=1}^{\infty}
2^k\mathbb{P}(X=2^k)\\
&=
\sum_{k=1}^{\infty}
2^k2^{-k}\\
&=
\sum_{k=1}^{\infty}1\\
&=\infty.
\end{aligned}
\]

Therefore,
\[
\boxed{
\mathbb{E}(X)=\infty.
}
\]

Thus the usual fair-price prescription
\(
\text{price}=\mathbb{E}(X)
\)
does not give a finite price.

\begin{observation}

The St.\ Petersburg game illustrates the difficulty caused by a
heavy-tailed distribution.

Although extremely large winnings are rare, they are sufficiently
large to make the expectation infinite.

\end{observation}


%============================================================
% Two Ways to Modify the Game
%============================================================

\subsection{Two Ways to Make the Game Finite}

There are at least two natural ways to modify the game.

\begin{enumerate}

    \item \textbf{Cap the winnings.}

    Impose an upper bound on the amount that can be won.

    \item \textbf{Limit the number of plays.}

    Instead of considering infinitely many independent plays, consider
    only a finite number of plays and determine an appropriate price
    for them.

\end{enumerate}

We now consider the second approach.


%============================================================
% Pricing n Independent Plays
%============================================================

\section{Pricing \(n\) Independent St.\ Petersburg Games}

\subsection{The Total Winning}

Let
\[
X_{n,1},X_{n,2},\ldots,X_{n,n}
\]
be independent copies of the St.\ Petersburg random variable $X$.

Thus
\[
X_{n,k}\overset{d}{=}X,
\qquad
1\leq k\leq n.
\]

Define
\[
S_n
=
\sum_{k=1}^nX_{n,k}.
\]

We want to understand the correct scale of $S_n$.

\begin{question}{Pricing \(n\) Independent St.\ Petersburg Games}
{pricing-n-st-petersburg-games}

How should one price $n$ independent plays of the St.\ Petersburg game?

\end{question}


%============================================================
% Weak Law for the Triangular Array
%============================================================

\subsection{Applying the Weak Law for Triangular Arrays}

We use the Weak Law for triangular arrays.

We seek a sequence $b_n\to\infty$ such that

\begin{align}
n\mathbb{P}(X>b_n)
&\longrightarrow0,
\tag{1}
\\
b_n^{-2}
n\mathbb{E}
\left[
X^2\mathbf{1}_{\{X\leq b_n\}}
\right]
&\longrightarrow0.
\tag{2}
\end{align}

These are precisely the two conditions needed to apply the triangular
array weak law with
\(
X_{n,k}=X_{k}
\)
and
\(
b_n
\)
as the normalization.

\subsection{The Tail Probability}

Let
\[
L_n
=
\left\lfloor
\log_2 b_n
\right\rfloor.
\]

Then
\[
2^{L_n}\leq b_n<2^{L_n+1}.
\]

Therefore,
\[
X>b_n
\]
implies
\[
X=2^k
\]
with
\[
k>L_n.
\]

Hence
\[
\begin{aligned}
\mathbb{P}(X>b_n)
&=
\sum_{k=L_n+1}^{\infty}2^{-k}\\
&=
2^{-L_n}.
\end{aligned}
\]

Since
\[
2^{L_n}\leq b_n<2^{L_n+1},
\]
we obtain
\[
\boxed{
\frac1{b_n}
<
\mathbb{P}(X>b_n)
\leq
\frac{2}{b_n}.
}
\tag{3}
\]

\subsection{Choice of \(b_n\)}

Choose
\[
\boxed{
b_n=n2^{k(n)},
}
\]
where
\[
k(n)\longrightarrow\infty.
\]

Then, by (3),
\[
\mathbb{P}(X>b_n)
\leq
\frac{2}{b_n}.
\]

Therefore,
\[
n\mathbb{P}(X>b_n)
\leq
\frac{2n}{b_n}
=
\frac{2}{2^{k(n)}}
\longrightarrow0.
\]

Thus condition (1) is satisfied.

\subsection{The Truncated Second Moment}

Let
\[
L_n
=
\left\lfloor
\log_2 b_n
\right\rfloor.
\]

Then
\[
\begin{aligned}
\mathbb{E}
\left[
X^2\mathbf{1}_{\{X\leq b_n\}}
\right]
&=
\sum_{k=1}^{L_n}
(2^k)^2\mathbb{P}(X=2^k)\\
&=
\sum_{k=1}^{L_n}
2^{2k}2^{-k}\\
&=
\sum_{k=1}^{L_n}2^k\\
&=
2^{L_n+1}-2.
\end{aligned}
\]

Since
\[
2^{L_n}\leq b_n,
\]
we have
\[
2^{L_n+1}\leq2b_n.
\]

Therefore,
\[
\mathbb{E}
\left[
X^2\mathbf{1}_{\{X\leq b_n\}}
\right]
\leq
2b_n.
\]

Consequently,
\[
\begin{aligned}
b_n^{-2}n
\mathbb{E}
\left[
X^2\mathbf{1}_{\{X\leq b_n\}}
\right]
&\leq
\frac{2n}{b_n}\\
&=
\frac{2}{2^{k(n)}}\\
&\longrightarrow0.
\end{aligned}
\]

Thus condition (2) is also satisfied.

\subsection{Conclusion from the Weak Law}

The Weak Law for triangular arrays therefore gives
\[
\frac{S_n-a_n}{b_n}
\xrightarrow{P}0,
\]
where
\[
a_n
=
n\mathbb{E}
\left[
X\mathbf{1}_{\{X\leq b_n\}}
\right].
\]

Now
\[
\begin{aligned}
\mathbb{E}
\left[
X\mathbf{1}_{\{X\leq b_n\}}
\right]
&=
\sum_{k=1}^{L_n}
2^k2^{-k}\\
&=
\sum_{k=1}^{L_n}1\\
&=
L_n.
\end{aligned}
\]

Hence
\[
a_n
=
nL_n
=
n
\left\lfloor
\log_2 b_n
\right\rfloor.
\]

Therefore,
\[
\boxed{
\frac{
S_n-
n\left\lfloor\log_2b_n\right\rfloor
}{b_n}
\xrightarrow{P}0.
}
\tag{4}
\]


%============================================================
% Finding the Natural Price
%============================================================

\section{The Natural Scale of the St.\ Petersburg Game}

\subsection{Choosing \(b_n\) for \(S_n/b_n\to1\)}

Ideally, we would like to choose $b_n$ such that
\[
\boxed{
\frac{S_n}{b_n}
\xrightarrow{P}1.
}
\]

Equivalently, for every $\varepsilon>0$,
\[
\mathbb{P}
\left(
\left|
\frac{S_n}{b_n}-1
\right|
>\varepsilon
\right)
\longrightarrow0.
\]

Using (4), it is enough to choose $b_n$ such that
\[
\frac{
n\left\lfloor\log_2b_n\right\rfloor
}{b_n}
\longrightarrow1.
\]

Thus we require
\[
\boxed{
\frac{
n\left\lfloor\log_2b_n\right\rfloor
}{b_n}
\longrightarrow1.
}
\tag{5}
\]

Now take
\[
b_n=n2^{k(n)}.
\]

Then
\[
\frac{
n\left\lfloor\log_2b_n\right\rfloor
}{b_n}
=
\frac{
\left\lfloor
\log_2n+k(n)
\right\rfloor
}{
2^{k(n)}
}.
\]

Therefore, we want
\[
\frac{
\log_2n+k(n)
}{
2^{k(n)}
}
\longrightarrow1.
\]

A natural choice is
\[
\boxed{
k(n)=\log_2\log_2 n.
}
\]

Indeed,
\[
2^{k(n)}
=
\log_2 n,
\]
and therefore
\[
\frac{
\log_2n+\log_2\log_2n
}{
\log_2n
}
\longrightarrow1.
\]

Thus we may choose
\[
\boxed{
b_n
=
n\log_2n.
}
\]

Consequently,
\[
\boxed{
\frac{S_n}{n\log_2n}
\xrightarrow{P}1.
}
\]

\begin{observation}

The natural scale of the total winnings in $n$ independent
St.\ Petersburg games is therefore
\[
n\log_2n.
\]

Hence the natural price per game grows like
\[
\boxed{
\log_2n.
}
\]

This is very different from the ordinary finite-expectation case,
where the natural price is a fixed number $\mathbb{E}(X)$.

\end{observation}


%============================================================
% Almost Sure Convergence
%============================================================

\section{From Convergence in Probability to Almost Sure Convergence}

\subsection{A Bernoulli Example}

Let
\[
X_1,X_2,\ldots
\]
be i.i.d. Bernoulli random variables with
\[
X_i\sim\operatorname{Ber}\left(\frac12\right).
\]

Define
\[
S_n=X_1+\cdots+X_n.
\]

The Weak Law gives
\[
\frac{S_n}{n}
\xrightarrow{P}
\frac12.
\]

We now ask whether
\[
\frac{S_n}{n}
\longrightarrow
\frac12
\]
almost surely.

\begin{question}{Almost Sure Convergence}
{almost-sure-convergence-bernoulli}

Does
\[
\frac{S_n}{n}
\longrightarrow
\frac12
\]
almost surely?

\end{question}


\subsection{Exponential Probability Bounds}

For every $\varepsilon>0$, Hoeffding's inequality gives
\[
\mathbb{P}
\left(
\left|
\frac{S_n}{n}-\frac12
\right|
>\varepsilon
\right)
\leq
2e^{-2n\varepsilon^2}.
\]

In particular, there exists a constant $c>0$ such that
\[
\boxed{
\mathbb{P}
\left(
\left|
\frac{S_n}{n}-\frac12
\right|
>\varepsilon
\right)
\leq
2e^{-cn}.
}
\tag{6}
\]

\subsection{Borel--Cantelli Argument}

Define
\[
A_n
=
\left\{
\left|
\frac{S_n}{n}-\frac12
\right|
>\varepsilon
\right\}.
\]

From (6),
\[
\mathbb{P}(A_n)
\leq
2e^{-cn}.
\]

Therefore,
\[
\sum_{n=1}^{\infty}\mathbb{P}(A_n)
\leq
2\sum_{n=1}^{\infty}e^{-cn}
<\infty.
\]

Hence, by the first Borel--Cantelli lemma,
\[
\mathbb{P}(A_n\text{ infinitely often})=0.
\]

Equivalently,
\[
\mathbb{P}
\left(
\left|
\frac{S_n}{n}-\frac12
\right|
>\varepsilon
\text{ infinitely often}
\right)
=0.
\]

Thus, for every $\varepsilon>0$,
\[
\left|
\frac{S_n}{n}-\frac12
\right|
\leq\varepsilon
\]
eventually almost surely.

Since this holds for every $\varepsilon>0$,
\[
\boxed{
\frac{S_n}{n}
\longrightarrow
\frac12
\qquad\text{almost surely}.
}
\]

\begin{observation}

This illustrates an important principle:

\[
\text{exponentially small deviation probabilities}
\quad+\quad
\text{Borel--Cantelli}
\]

\[\quad\Longrightarrow\quad
\text{almost sure convergence}.
\]

\end{observation}


%============================================================
% Contradiction Formulation
%============================================================

\subsection{Equivalent Contradiction Argument}

Suppose, for contradiction, that
\[
\frac{S_n}{n}
\not\longrightarrow
\frac12
\qquad\text{almost surely}.
\]

Then
\[
\mathbb{P}
\left(
\frac{S_n}{n}
\not\longrightarrow
\frac12
\right)
>0.
\]

Equivalently, there exist
\[
\theta>0
\qquad\text{and}\qquad
\varepsilon>0
\]
such that
\[
\mathbb{P}
\left(
\left|
\frac{S_n}{n}-\frac12
\right|
>\varepsilon
\text{ infinitely often}
\right)
>\theta.
\]

But the Borel--Cantelli argument above gives
\[
\mathbb{P}
\left(
\left|
\frac{S_n}{n}-\frac12
\right|
>\varepsilon
\text{ infinitely often}
\right)
=0,
\]
which is a contradiction.

Therefore,
\[
\boxed{
\frac{S_n}{n}
\longrightarrow
\frac12
\qquad\text{almost surely}.
}
\]


%============================================================
% Summary
%============================================================



%============================================================
% End of Lecture
%============================================================

\lectureend
